\section*{C6: Rezoluție}

\subsection{Exemplu de respingere prin rezoluție}

Pentru un scop
\[
	G=\forall(\neg A_1\lor\cdots\lor\neg A_m)
\]
și o clauză redenumită
\[
	C=\forall(B_0\lor\neg B_1\lor\cdots\lor\neg B_k),
\]
alegem $A_i$ și $B_0$ cu același simbol de relație. Dacă $\theta$ este
un cgu al lor, noul scop se obține înlocuind $A_i$ cu corpul clauzei și
aplicând $\theta$.

Considerăm programul
\[
	P=\bigl\{
		p(a),\quad
		p(f(x))\leftarrow p(x)
	\bigr\}
\]
și scopul
\[
	G_0=\forall\neg p(f(f(y))).
\]
La fiecare pas redenumim variabila clauzei recursive:
\[
	\begin{aligned}
	G_0
	&\xRightarrow[\theta_0=\{x_1\mapsto f(y)\}]
	{\,p(f(x_1))\leftarrow p(x_1)\,}
	G_1=\forall\neg p(f(y)),\\
	G_1
	&\xRightarrow[\theta_1=\{x_2\mapsto y\}]
	{\,p(f(x_2))\leftarrow p(x_2)\,}
	G_2=\forall\neg p(y),\\
	G_2
	&\xRightarrow[\theta_2=\{y\mapsto a\}]
	{\,p(a)\,}
	G_3=\bot.
	\end{aligned}
\]
Aceasta este o $P$-respingere a lui $G_0$. Substituția calculată,
restricționată la variabilele scopului inițial, este
\[
	(\widetilde{\theta}_2\circ\widetilde{\theta}_1\circ\theta_0)
	\big|_{\{y\}}=\{y\mapsto a\}.
\]
Prin teorema de corectitudine,
\[
	P\models p(f(f(a))).
\]
